%=====================================================================
\chapter{Greedy Algorithms}\label{ch:greedy}
%=====================================================================
\locator{4}{Part IV --- The Design Paradigms}

\begin{outcomes}
By the end of this chapter you will be able to:
\begin{tight}
  \item \textbf{Define} the greedy strategy and \textbf{apply} it to a stated
        problem \emph{(CLO-6)}.
  \item \textbf{State} the greedy-choice property and optimal substructure, and
        \textbf{explain} why both are needed.
  \item \textbf{Prove} a greedy algorithm optimal by an exchange argument.
  \item \textbf{Construct} a counter-example showing a plausible greedy rule is
        not optimal.
  \item \textbf{Describe} Huffman coding and \textbf{relate} its output to the
        entropy bound.
\end{tight}
\end{outcomes}

\begin{prereq}
\chref{bruteforce} for the baseline and the oracle technique. Nothing else in
Part~II is required, though \chref{structures}'s heap is the data structure
Huffman's algorithm runs on.
\end{prereq}

\begin{keyterms}
greedy algorithm $\bullet$ greedy-choice property $\bullet$ optimal
substructure $\bullet$ exchange argument $\bullet$ activity selection
$\bullet$ counter-example $\bullet$ Huffman code $\bullet$ prefix code
$\bullet$ entropy $\bullet$ matroid
\end{keyterms}

\section{Commit and Never Reconsider}

\begin{definitionbox}[Greedy Algorithm]
A \term{greedy algorithm} builds a solution one step at a time, at each step
taking the option that looks best \emph{right now}, by some local rule, and
never revisiting that decision.

\smallskip
It is the simplest strategy in the book and the one most often wrong. What
makes it interesting is that when it works, it is usually both the fastest
algorithm available and the easiest to implement --- and whether it works is
decidable by a proof, not by testing.
\end{definitionbox}

\begin{conceptbox}[The two properties]
A greedy algorithm is optimal when the problem has both of these.

\smallskip
\textbf{1. The greedy-choice property.} A globally optimal solution can be
reached by making a locally optimal choice. More precisely: there exists an
optimal solution that \emph{contains} the greedy first choice. Note the
``there exists'' --- greedy need not find \emph{the} optimal solution, only
\emph{an} optimal one.

\smallskip
\textbf{2. Optimal substructure.} An optimal solution to the problem contains
within it optimal solutions to the subproblems that remain after the choice.

\smallskip
\textbf{Both are required and they are different.} Optimal substructure alone
gives you dynamic programming (\chref{dpone}); it is the greedy-choice property
that lets you commit without exploring alternatives, and it is the one that
usually fails. \chref{dpone} returns to this comparison in full.
\end{conceptbox}

\section{Activity Selection}

\begin{definitionbox}[Activity Selection]
Given $n$ activities, activity $i$ occupying the half-open interval
$[s_{i}, f_{i})$, choose a largest set of mutually compatible activities --- no
two overlapping.
\end{definitionbox}

Three local rules suggest themselves, and they sound equally sensible:

\begin{tight}
  \item take the activity that \textbf{finishes earliest};
  \item take the \textbf{shortest} activity;
  \item take the activity that \textbf{starts earliest}.
\end{tight}

\dsfig{\aoaalg{\algname{Activity-Select}$(s, f, n)$}{
\State sort the activities so that $f_{1} \le f_{2} \le \cdots \le f_{n}$
\State $A \gets \{1\}$;\quad $k \gets 1$
  \Comment{the first to finish is always taken}
\For{$m \gets 2$ \textbf{to} $n$}
  \If{$s_{m} \ge f_{k}$}
    \State $A \gets A \cup \{m\}$
      \Comment{compatible: commit, never reconsider}
    \State $k \gets m$
  \EndIf
\EndFor
\State \Return $A$
}}{\algname{Activity-Select}, greedy by earliest finishing time. After the
sort it is a single pass, so the whole algorithm is
$\bigTh{n\lg n}$.}{activityselect}

\begin{theorem}[Earliest finish is optimal]\label{thm:activity}
Let $S$ be a set of activities and let $a_{1}$ be one with the earliest
finishing time. Then some maximum-size set of mutually compatible activities
contains $a_{1}$.
\end{theorem}

\begin{proof}
This is an \term{exchange argument}: take any optimal solution and show it can
be modified to contain the greedy choice without getting worse.

\smallskip
Let $A$ be a maximum-size compatible set, and let $a_{j}$ be the activity in
$A$ with the earliest finishing time. If $a_{j} = a_{1}$ there is nothing to
prove. Otherwise consider
\[ A' = \bigl(A \setminus \{a_{j}\}\bigr) \cup \{a_{1}\}. \]

\smallskip
\textbf{$A'$ is compatible.} The activities of $A$ other than $a_{j}$ all start
at or after $f_{j}$, because $a_{j}$ finishes earliest in $A$ and $A$ is
compatible. Since $a_{1}$ has the earliest finishing time in all of $S$, we
have $f_{1} \le f_{j}$, so every other activity in $A$ also starts at or after
$f_{1}$ and does not overlap $a_{1}$.

\smallskip
\textbf{$A'$ has the same size as $A$}, since one element was exchanged for
one. So $A'$ is also a maximum-size compatible set, and it contains $a_{1}$.

\smallskip
Optimal substructure then completes the induction: after committing to
$a_{1}$, the remaining problem is activity selection on
$\{a_{i} : s_{i} \ge f_{1}\}$, an instance of the same problem, and the same
argument applies to it.
\end{proof}

\begin{worked}[Measured: the Three Rules, Against an Oracle]
Rules 2 and 3 have no such proof. Rather than assert they are wrong,
\texttt{code/ch14\_greedy.cpp} checks all three against exhaustive search
(\chref{bruteforce}'s oracle) on 20{,}000 random instances of 10 activities:

\rowcolors{2}{white}{lightbg}
\begin{center}
\begin{tabular}{@{}lr@{}}
\toprule
\rowcolor{primary!15}
greedy rule & instances solved optimally \\
\midrule
earliest \textbf{finish} time & \textbf{100.0\%} \\
\textbf{shortest} activity    & 94.5\% \\
earliest \textbf{start} time  & 62.1\% \\
\bottomrule
\end{tabular}
\end{center}

\smallskip
\textbf{Read the middle row carefully.} ``Take the shortest'' is right 94.5\%
of the time. Any amount of testing on random inputs would leave you confident
in it, and it is wrong. A 94.5\% algorithm is not a 94.5\%-correct algorithm;
it is an incorrect algorithm that is usually lucky.

\smallskip
\textbf{Read the last row too.} ``Take whatever starts first'' --- the most
natural rule of all, and the one a scheduler written in a hurry implements ---
is optimal barely three-fifths of the time.

\smallskip
\textbf{And the top row is 100.0\%, not 100\%-ish}, because it is a theorem.
That is the difference between a proof and a benchmark, and it is the reason
Theorem~\ref{thm:activity} is in this chapter.
\end{worked}

\section{When Greedy Fails: Coin Change}

\begin{worked}[The Same Algorithm, Two Coin Systems]
Make an amount from coins of given denominations, using as few coins as
possible. Greedy: take the largest coin that fits, repeatedly.

\rowcolors{2}{white}{lightbg}
\begin{center}
\begin{tabular}{@{}lrl@{}}
\toprule
\rowcolor{primary!15}
coin system & greedy wrong on & first failure \\
\midrule
$\{1, 5, 10, 25\}$ & 0 of 200 amounts & --- \\
$\{1, 3, 4\}$ & \textbf{49 of 200 amounts} & amount 6 \\
\bottomrule
\end{tabular}
\end{center}

\smallskip
\textbf{The counter-example is six.} Greedy takes 4, then 1, then 1 --- three
coins. The optimum is $3 + 3$ --- two coins. Taking the largest coin was the
mistake, and no amount of care in the implementation can fix a wrong strategy.

\smallskip
\textbf{The lesson is not that greedy is unreliable.} It is that greediness is
a property of the \emph{problem instance}, not of the approach. The identical
algorithm is provably optimal for $\{1,5,10,25\}$ --- which is why it works for
every currency you have handled --- and provably wrong for $\{1,3,4\}$. You
cannot tell which case you are in by looking at the code.

\smallskip
\textbf{And when greedy fails, dynamic programming does not.} \chref{dpone}'s
method solves coin change optimally for any denomination set, in
$\bigTh{nW}$ --- slower, and always right.
\end{worked}

\begin{alertbox}[How to approach a greedy algorithm]
\begin{tightnum}
  \item \textbf{Write the greedy rule down precisely.} ``Greedy'' names a
        family; the three activity-selection rules are three different
        algorithms with three different success rates.
  \item \textbf{Try to prove the greedy-choice property} by exchange argument.
        If the proof goes through, you are finished.
  \item \textbf{If it does not, look for the counter-example} where the proof
        stuck. Counter-examples for greedy algorithms are usually tiny ---
        $\{1,3,4\}$ and the amount 6.
  \item \textbf{Never conclude from testing.} 94.5\% looked like success.
\end{tightnum}
\end{alertbox}

\section{Huffman Coding}

\begin{definitionbox}[Prefix Code]
A \term{prefix code} assigns each symbol a binary string such that no
codeword is a prefix of another. That property is what makes a code
unambiguously decodable without separators: reading left to right, the first
codeword you complete is the right one.

\smallskip
A prefix code corresponds exactly to a binary tree whose leaves are the
symbols; the codeword is the path, with 0 for left and 1 for right. The cost of
a tree $T$ is
\[ B(T) = \sum_{c} f(c)\, d_{T}(c), \]
the total encoded length, where $f(c)$ is $c$'s frequency and $d_{T}(c)$ its
depth.
\end{definitionbox}

\dsfig{\aoaalg{\algname{Huffman}$(C)$}{
\State $Q \gets$ a min-priority queue of all symbols, keyed by frequency
\For{$i \gets 1$ \textbf{to} $|C| - 1$}
  \State allocate a new node $z$
  \State $z.\mathit{left} \gets$ \algname{Extract-Min}$(Q)$
  \State $z.\mathit{right} \gets$ \algname{Extract-Min}$(Q)$
    \Comment{the two rarest}
  \State $z.\mathit{freq} \gets z.\mathit{left}.\mathit{freq}
         + z.\mathit{right}.\mathit{freq}$
  \State \algname{Insert}$(Q, z)$
\EndFor
\State \Return \algname{Extract-Min}$(Q)$
  \Comment{the root}
}}{\algname{Huffman}. The greedy choice is on line 5: the two \emph{least}
frequent symbols are merged first, which puts them deepest. With a binary heap
the running time is $\bigTh{n \lg n}$.}{huffman}

\begin{lemma}[The greedy choice is safe]\label{lem:huffman}
Let $x$ and $y$ be two symbols of lowest frequency in $C$. Then there is an
optimal prefix code for $C$ in which $x$ and $y$ have the same (maximal) depth
and differ only in the last bit.
\end{lemma}

\begin{proof}
Let $T$ be an optimal tree, and let $a$ and $b$ be two sibling leaves of
maximum depth in $T$ (they exist: the deepest leaf has a sibling, or the tree
could be improved by promoting it). Assume without loss of generality
$f(a) \le f(b)$ and $f(x) \le f(y)$.

\smallskip
Exchange $x$ with $a$ to form $T'$. The change in cost is
\[ B(T) - B(T') = \bigl(f(a) - f(x)\bigr)\bigl(d_{T}(a) - d_{T}(x)\bigr)
   \;\ge\; 0, \]
because $f(x) \le f(a)$ ($x$ is of lowest frequency) and
$d_{T}(x) \le d_{T}(a)$ ($a$ is of maximal depth), making both factors
non-negative. So $B(T') \le B(T)$. Exchanging $y$ with $b$ in $T'$ gives $T''$
with $B(T'') \le B(T') \le B(T)$ by the same argument.

\smallskip
Since $T$ was optimal, $B(T'') = B(T)$, so $T''$ is also optimal --- and in it
$x$ and $y$ are sibling leaves at maximum depth.
\end{proof}

\begin{theorem}[Huffman's algorithm is optimal]\label{thm:huffman}
\algname{Huffman} produces a prefix code of minimum total length.
\end{theorem}

\begin{proof}
By induction on $|C|$. The base case $|C| = 2$ is immediate: both codewords
have length 1 and no code can do better.

\smallskip
For the inductive step, let $x, y$ be the two least frequent symbols and let
$C'$ be $C$ with $x$ and $y$ replaced by a single symbol $z$ of frequency
$f(x)+f(y)$. By the inductive hypothesis, \algname{Huffman} produces an optimal
tree $T'$ for $C'$. Replacing $z$'s leaf by an internal node with children $x$
and $y$ gives a tree $T$ for $C$ with
\[ B(T) = B(T') + f(x) + f(y), \]
since $x$ and $y$ sit one level below where $z$ was. The $+f(x)+f(y)$ is a
constant independent of the tree, so minimising $B(T')$ minimises $B(T)$ over
all trees in which $x$ and $y$ are siblings --- and Lemma~\ref{lem:huffman}
says some optimal tree is of that form. Hence $T$ is optimal.
\end{proof}

\begin{worked}[Huffman Against the Entropy Bound]
Shannon's source coding theorem says no prefix code can average fewer than
$H = -\sum_{c} p(c)\lg p(c)$ bits per symbol, and Huffman is always within one
bit of that. Both halves are checkable.

\rowcolors{2}{white}{lightbg}
\begin{center}
\begin{tabular}{@{}lrrrr@{}}
\toprule
\rowcolor{primary!15}
text & symbols & entropy $H$ & Huffman & saving vs 8-bit \\
\midrule
English prose     &  22 & 3.9670 & 4.0130 & 49.8\% \\
DNA (4 symbols)   &   4 & 1.9999 & 2.0000 & 75.0\% \\
95\% one symbol   &   6 & 0.4065 & \textbf{1.1202} & 86.0\% \\
uniform bytes     & 256 & 7.9953 & 8.0000 &  0.0\% \\
\bottomrule
\end{tabular}
\end{center}

\smallskip
\textbf{Every row satisfies $H \le \text{Huffman} < H+1$}, as the theory
requires. Three of the four are within 0.05 bits of the floor.

\smallskip
\textbf{The third row is the interesting one.} Entropy 0.4065, Huffman 1.1202
--- a gap of 0.71 bits, nearly the whole slack the theorem allows. The reason
is structural: Huffman assigns each symbol a \emph{whole number} of bits, and
the minimum is 1. When one symbol has probability 0.95 its information content
is $-\lg 0.95 = 0.074$ bits, and Huffman must spend 1. It cannot do better,
and the loss is a factor of 2.8.

\smallskip
\textbf{That is the known limit of Huffman coding}, and the reason arithmetic
coding exists --- it can spend fractional bits per symbol and approaches $H$
arbitrarily closely. A greedy algorithm that is provably optimal \emph{within
its class} can still be far from the information-theoretic bound, because the
class itself is the constraint. Compare \chref{lowerbound}, where merge sort
came within 1\% of its floor; here the restriction to whole bits costs 175\%.

\smallskip
\textbf{And the last row.} With 256 equally likely bytes, $H = 8$ and Huffman
produces exactly 8 bits per symbol --- it saves nothing, correctly. Compression
is possible only where there is structure to exploit, and a greedy algorithm
cannot create any.
\end{worked}

\begin{pitfall}
\begin{tight}
  \item \textbf{Saying ``use a greedy algorithm'' without naming the rule.}
        Three rules for activity selection scored 100\%, 94.5\% and 62.1\%.
  \item \textbf{Concluding optimality from testing.} ``Take the shortest'' was
        right in 94.5\% of 20{,}000 random trials and is wrong.
  \item \textbf{Assuming greedy works because it works for coins.} It does ---
        for $\{1,5,10,25\}$. For $\{1,3,4\}$ it fails at the amount 6.
  \item \textbf{Confusing the two properties.} Optimal substructure alone gives
        dynamic programming. It is the greedy-choice property that permits
        committing without exploring.
  \item \textbf{Proving the wrong direction in an exchange argument.} You must
        show the \emph{optimal} solution can be modified to contain the greedy
        choice --- not that the greedy solution can be modified.
  \item \textbf{Expecting Huffman to reach the entropy.} It cannot spend
        fractional bits; measured, it was 175\% above $H$ on a highly skewed
        source.
  \item \textbf{Using Huffman on a uniform source.} Measured saving: 0.0\%.
\end{tight}
\end{pitfall}

\begin{lab}[Three Rules, Two Coin Systems, One Code]
Reproduces all three tables. Expect twenty-five minutes; part A runs 20{,}000
brute-force searches.

\begin{lstlisting}[language=C++]
// ch14_greedy.cpp -- when a local choice can be trusted, and when not.
// Build: cl /O2 /EHsc /std:c++17 ch14_greedy.cpp     (see Appendix A)

// --- 1. Activity selection, greedy by earliest finish --------------
// Sort by finishing time and take anything compatible. The claim is
// that the activity finishing first is in SOME optimal solution --
// the greedy-choice property, and what the chapter proves.
static int greedy_finish(std::vector<Act> a) {
    std::sort(a.begin(), a.end(),
              [](const Act& x, const Act& y) {
                  return x.finish < y.finish; });
    int count = 0, last = -1;
    for (const Act& x : a)
        if (x.start >= last) { ++count; last = x.finish; }
    return count;
}

// --- 2. Two plausible rules that are simply wrong ------------------
// Both "take the shortest" and "take the earliest starting" sound
// just as reasonable as rule 1. Neither is optimal, and the program
// finds the counter-examples by brute force rather than asserting it.

// --- 3. The oracle: every subset, checked --------------------------
// Chapter 7's lesson. n <= 18 so 2^n is affordable, and it is the
// only way to be sure the greedy answers are right, not plausible.

// --- 4. Coin change: greedy, and the optimum it sometimes misses ---
// Greedy takes the largest coin that fits. For {1, 5, 10, 25} that is
// optimal; for {1, 3, 4} it is not, and the smallest counter-example
// is six. Dynamic programming (Chapter 15) gets it right for any
// system.
static int coins_greedy(std::vector<int> c, int amount) {
    std::sort(c.rbegin(), c.rend());
    int used = 0;
    for (int v : c) while (amount >= v) { amount -= v; ++used; }
    return amount == 0 ? used : -1;
}

// --- 5. Huffman coding ---------------------------------------------
// Repeatedly merge the two least frequent symbols. The greedy choice
// is that the two rarest symbols can be put deepest, as siblings.
// Total encoded length = sum over leaves of freq * depth, which is
// also the sum of every internal node's frequency -- a neater way to
// compute it than walking the tree.

// --- 6. Is the greedy rule actually optimal? -----------------------
// --- 7. Coin change: the same algorithm, two coin systems ----------
// --- 8. Huffman against the entropy bound --------------------------
// Shannon: no prefix code can average fewer than H bits per symbol,
// and Huffman is always within 1 bit of it. Both are checkable.
\end{lstlisting}

\textbf{What to notice.} Part 6's middle column. A rule that is right 94.5\% of
the time is not 94.5\% of a correct algorithm --- it is wrong, and the 5.5\% is
where it will meet your data.

\smallskip
\textbf{Then try to break it.} Find the smallest instance on which ``take the
shortest'' fails, by searching exhaustively over instances rather than
sampling. Three activities suffice: $[0,10)$, $[9,12)$ and $[11,20)$ --- the
shortest is the middle one, taking it forbids both others, and the optimum
takes the outer two. Having the counter-example in hand is worth more than the
94.5\%.
\end{lab}

\begin{checkpoint}
\begin{tightnum}
  \item State the greedy-choice property and optimal substructure, and give a
        problem that has the second but not the first.
  \item Greedy coin change fails for $\{1,3,4\}$ at the amount 6. Find the next
        amount at which it fails, and say why the counter-examples are small.
  \item Huffman achieved 1.1202 bits per symbol against an entropy of 0.4065.
        Explain the gap without blaming the algorithm.
\end{tightnum}
\end{checkpoint}

\begin{chaptersummary}
\begin{tight}
  \item A greedy algorithm commits to the locally best choice and never
        reconsiders. It is optimal when the problem has both the greedy-choice
        property and optimal substructure.
  \item Optimality is proved by an \emph{exchange argument}: show that some
        optimal solution can be modified to contain the greedy choice without
        getting worse.
  \item Measured: of three plausible activity-selection rules, earliest-finish
        was optimal on 100.0\% of 20{,}000 random instances, shortest-first on
        94.5\% and earliest-start on 62.1\%. Testing cannot distinguish a
        theorem from luck.
  \item Measured: greedy coin change is optimal for $\{1,5,10,25\}$ and wrong
        on 49 of 200 amounts for $\{1,3,4\}$, first failing at 6. Greediness is
        a property of the instance, not of the code.
  \item Huffman's algorithm is greedy --- merge the two rarest symbols --- and
        provably produces a minimum-length prefix code.
  \item Measured: Huffman lay in $[H, H+1)$ on all four corpora. On a 95\%-skewed
        source it used 1.1202 bits against an entropy of 0.4065, because it
        cannot spend fractional bits; on a uniform source it saved nothing,
        correctly.
\end{tight}
\end{chaptersummary}

\begin{reviewq}
  \item \textbf{(MCQ)} The greedy-choice property states that:
        (a)~the greedy solution is unique; (b)~some optimal solution contains
        the greedy first choice; (c)~every optimal solution is greedy;
        (d)~the problem has no optimal substructure.
  \item \textbf{(MCQ)} Huffman coding produces a code whose average length is:
        (a)~always equal to the entropy; (b)~at least the entropy and less than
        entropy $+1$; (c)~always 8 bits; (d)~unrelated to the entropy.
  \item \textbf{(Short)} Define the greedy strategy and apply it to activity
        selection, stating the rule precisely \emph{(CLO-6)}.
  \item \textbf{(Short)} Explain an exchange argument, and say which direction
        the modification goes.
  \item \textbf{(Short)} Why is a greedy rule that succeeds on 94.5\% of random
        instances not acceptable?
  \item \textbf{(Applied)} Prove that greedy coin change is optimal for
        $\{1,5,10,25\}$, or give the argument's outline and identify where it
        would fail for $\{1,3,4\}$.
  \item \textbf{(Applied)} You must schedule jobs with deadlines and profits,
        one machine, each job taking unit time, to maximise profit. Propose a
        greedy rule, attempt an exchange argument, and either prove it or give
        a counter-example.
\end{reviewq}
